\section{Considerações iniciais}
\label{sec:consideracoes-iniciais-referencial}

Este capítulo apresenta os fundamentos conceituais, normativos e educacionais necessários para compreender a relação entre a BNCC Computação e a formação de professores. Inicialmente, o capítulo delimita os conhecimentos, as práticas e as formas de raciocínio da Computação que constituem objeto de ensino e aprendizagem na Educação Básica. Em seguida, examina como esses conhecimentos foram incorporados à política curricular brasileira, as responsabilidades decorrentes de sua implementação e as indefinições normativas relacionadas à formação docente.

Na sequência, são discutidos os desafios curriculares, formativos, materiais e institucionais da implementação do ensino de Computação. O capítulo analisa o papel da formação docente, os limites de uma resposta baseada apenas na formação continuada, as possibilidades da formação inicial e os conhecimentos necessários à preparação para o ensino da área. Esses fundamentos sustentam a discussão desenvolvida no Capítulo~\ref{chapter:revisao-integracao-curricular}, que examina as diferentes licenciaturas, caracteriza a integração da Computação como problema de decisão curricular e delimita a lacuna apresentada na Seção~\ref{sec:lacuna}.

\section{A Computação como objeto de ensino e aprendizagem na Educação Básica}
\label{sec:computacao-conhecimento-escolar}

\subsection{Delimitação conceitual do ensino de Computação}
\label{subsec:delimitacao-educacao-computacao}

A literatura emprega diferentes expressões para tratar da presença da Computação nos processos educacionais. A expressão \textit{Educação em Computação} designa o campo mais amplo que investiga o ensino e a aprendizagem da Computação e abrange temas como currículo, práticas pedagógicas, avaliação, materiais didáticos, formação de professores e políticas educacionais. A expressão \textit{ensino de Computação}, por sua vez, refere-se às situações em que conhecimentos, práticas e formas de raciocínio próprios da Computação são tratados como conteúdos a serem ensinados e aprendidos.

O ensino de Computação será, portanto, o termo principal utilizado neste trabalho para designar o objeto formativo discutido. A expressão Educação em Computação será reservada às referências ao campo mais amplo de pesquisa. Essa distinção também ajuda a separar o ensino dos conhecimentos da Computação do uso de computadores, plataformas e outros recursos digitais para ensinar conteúdos de diferentes áreas.

Não existe uma nomenclatura internacional única para designar o ensino de Computação na Educação Básica. Termos como \textit{computer science}, \textit{computing} e \textit{informatics} assumem sentidos diferentes conforme a tradição científica e curricular de cada país \cite{heintz2016review}.

Essa distinção aparece, por exemplo, no relatório que a Royal Society, academia de ciências do Reino Unido, publicou em \citeyear{royal2012shut} sobre o ensino de computação nas escolas britânicas. Nesse relatório, a instituição propõe distinguir Ciência da Computação, Tecnologia da Informação e Letramento Digital \cite{royal2012shut}. A Ciência da Computação trata de princípios, conceitos e processos relacionados à informação e aos sistemas computacionais.

Essa distinção é necessária porque o uso de uma tecnologia em uma atividade educacional não significa, por si só, que conhecimentos de computação estejam sendo ensinados. Por exemplo, uma plataforma pode ser usada para apresentar conteúdos sem que os estudantes aprendam como informações são representadas, como algoritmos são construídos, como sistemas digitais funcionam ou quais impactos decorrem de seu desenvolvimento e uso.

\subsection{Referenciais para o ensino de Computação na Educação Básica}
\label{subsec:referenciais-computacao-escolar}

Referenciais curriculares elaborados em diferentes países tratam a Computação como um campo de conhecimentos, práticas e formas de raciocínio a ser ensinado na Educação Básica. Nos Estados Unidos, uma coalizão liderada pela \sigla{ACM}{Association for Computing Machinery}, pelo Code.org e pela \sigla{CSTA}{Computer Science Teachers Association}, publicou em 2016 o K--12 Computer Science Framework. O documento orienta estados e redes de ensino norte-americanas na elaboração de currículos e padrões para o ensino de Computação na Educação Básica. Ele organiza o campo em cinco conceitos principais: sistemas computacionais; redes e Internet; dados e análise; algoritmos e programação; e impactos da computação~\cite{k2016k}. Além dos conceitos, o documento propõe práticas como decompor problemas, desenvolver abstrações, criar e testar artefatos computacionais e comunicar ideias sobre computação.

O currículo de Tecnologia da Nova Zelândia \cite{NewZealandMoE2017}, revisado em 2017 pelo Ministério da Educação do país, também trata a computação como parte da formação escolar obrigatória. Dentro da área de tecnologia, o currículo criou duas frentes dedicadas às tecnologias digitais, Pensamento Computacional e Desenvolvimento de Produtos Digitais. A primeira trata da decomposição de problemas, da construção de algoritmos e da depuração de erros. A segunda trata da criação, do teste e da avaliação de produtos digitais voltados a usuários reais. Essas frentes se tornaram parte obrigatória do currículo do Ensino Fundamental neozelandês em 2020 \cite{NewZealandMoE2017}.

No Brasil, a \sigla{SBC}{Sociedade Brasileira de Computação} apresentou uma proposta de delimitação curricular nas \textit{Diretrizes para o Ensino de Computação na Educação Básica} \cite{ribeiro2019diretrizes}. O documento compreende a computação como uma área de conhecimento dedicada ao estudo de processos de informação e ao desenvolvimento de formas de representar, analisar e resolver problemas. Também diferencia o ensino de computação da fluência digital e do uso pedagógico de tecnologias. As diretrizes da SBC organizam os conhecimentos de Computação previstos para a Educação Básica em três eixos articulados: Pensamento Computacional, Mundo Digital e Cultura Digital. Também propõem uma progressão entre experiências concretas, construção de abstrações, programação, compreensão de sistemas computacionais e análise dos impactos da Computação \cite{ribeiro2019diretrizes}.

Apesar das diferenças de organização, os referenciais apresentados convergem ao reconhecer que a computação possui conceitos, práticas e formas próprias de compreender problemas. A presença desses conhecimentos na Educação Básica é justificada pela influência dos sistemas computacionais sobre a comunicação, o trabalho, a ciência, a economia e a participação social. Sua aprendizagem pode ajudar os estudantes a compreender os princípios, as possibilidades, as limitações e os impactos das tecnologias que utilizam \cite{caspersen2019informatics,royal2012shut}.

\subsection{Definição de ensino de Computação adotada}
\label{subsec:definicao-computacao-pesquisa}

Como esta pesquisa está situada na Educação Básica brasileira, a definição adotada toma como referências principais as diretrizes da Sociedade Brasileira de Computação \cite{ribeiro2019diretrizes} e a organização estabelecida pela BNCC Computação \cite{brasil2022bncccomputacao}.

Neste trabalho, a expressão \textit{ensino de Computação} designa o ensino e a aprendizagem dos conhecimentos, das práticas e das formas de raciocínio abrangidos pelos eixos Pensamento Computacional, Mundo Digital e Cultura Digital. O termo inclui, portanto, tanto a compreensão dos conceitos da área quanto o desenvolvimento de práticas de resolução de problemas, criação de soluções e análise dos sistemas computacionais e de seus impactos. Essa definição estabelece o objeto formativo discutido nas seções seguintes.

\section{A BNCC Computação e a demanda curricular nacional}
\label{sec:bncc-computacao}

%\subsection{Da BNCC ao complemento de Computação}
%\label{subsec:bncc-complemento-computacao}

A BNCC define as aprendizagens essenciais que os estudantes devem desenvolver ao longo da Educação Básica \cite{brasil2018bncc}. A BNCC da Educação Infantil e do Ensino Fundamental foi instituída pela Resolução CNE/CP nº 2/2017, enquanto a etapa do Ensino Médio foi instituída pela Resolução CNE/CP nº 4/2018 \cite{brasil2017bncc,brasil2018bnccem}.

A BNCC já tratava da cultura digital e do uso crítico, responsável e criativo das tecnologias \cite{brasil2018bncc}. Entretanto, os conhecimentos próprios da computação não apareciam de forma específica e progressiva. A própria Resolução CNE/CP nº 2/2017 previa a elaboração posterior de normas complementares sobre o tema \cite{brasil2017bncc}.

Em 2022, o Parecer CNE/CEB nº 2/2022 apresentou a proposta técnica para a inclusão da Computação na Educação Básica \cite{brasil2022parecercomputacao2}. Depois de homologado, o parecer deu origem à Resolução CNE/CEB nº 1/2022, que instituiu as Normas sobre Computação na Educação Básica em complemento à BNCC \cite{brasil2022resolucao1}. Anexas à resolução estão as tabelas de competências, objetos de conhecimento, objetivos de aprendizagem e habilidades para cada etapa da escolarização \cite{brasil2022bncccomputacao}.

Neste trabalho, a expressão BNCC Computação refere-se ao conjunto formado pela Resolução CNE/CEB nº 1/2022 e por suas tabelas anexas. Não se trata de uma base curricular independente, mas de um complemento à BNCC existente. Seu alcance compreende a Educação Infantil, os Anos Iniciais e Finais do Ensino Fundamental e o Ensino Médio.

\subsection{Os três eixos da BNCC Computação}
\label{subsec:eixos-bncc-computacao}

A BNCC Computação organiza as aprendizagens em três eixos que se relacionam entre si: Pensamento Computacional, Mundo Digital e Cultura Digital \cite{brasil2022bncccomputacao}.

O \textbf{Pensamento Computacional} trata da representação, da análise e da resolução de problemas com base em conceitos e práticas da Computação. Inclui, entre outros elementos, reconhecimento de padrões, decomposição, abstração, modelagem, construção de algoritmos, programação, teste e comparação de soluções.

O \textbf{Mundo Digital} aborda os elementos que constituem e sustentam as tecnologias digitais. Compreende temas como dispositivos, hardware, software, representação de dados, redes, Internet, protocolos, segurança e comunicação entre sistemas.

A \textbf{Cultura Digital} trata das relações entre Computação, sociedade e vida cotidiana. Abrange uso responsável das tecnologias, segurança, privacidade, autoria, participação em ambientes digitais, avaliação crítica de informações e impactos sociais, políticos, econômicos e ambientais.

Os três eixos não representam conteúdos isolados. A construção de um programa, por exemplo, pode mobilizar formas de resolução de problemas associadas ao Pensamento Computacional, depender da compreensão de elementos do Mundo Digital e envolver decisões de autoria, privacidade ou impacto social próprias da Cultura Digital.

O complemento também propõe uma progressão ao longo da Educação Básica. As experiências concretas e lúdicas da Educação Infantil dão lugar à introdução gradual de conceitos e práticas durante o Ensino Fundamental. No Ensino Médio, espera-se uma abordagem mais abstrata, crítica e aprofundada. Essa progressão é relevante para a formação docente, pois quem ensina precisa compreender não apenas os conteúdos de uma etapa, mas também suas relações com aprendizagens anteriores e posteriores.

\subsection{Implementação e limites da norma}
\label{subsec:implementacao-limites-norma}

A Resolução CNE/CEB nº 1/2022 não se limita a definir aprendizagens. Ela determina que os currículos escolares considerem as tabelas anexas e que a formação inicial e continuada de professores leve em conta os conhecimentos e as habilidades previstos no complemento \cite{brasil2022resolucao1}.

A responsabilidade por definir parâmetros e abordagens pedagógicas de implementação é atribuída aos Estados, Municípios e ao Distrito Federal. Ao Ministério da Educação cabem ações de formação, apoio curricular, elaboração de materiais, avaliação e assessoramento aos sistemas de ensino \cite{brasil2022resolucao1}.

Essa distribuição estabelece uma demanda nacional de implementação, mas não prescreve um único arranjo curricular para atendê-la. A resolução não determina que a Computação seja organizada como disciplina autônoma, não fixa carga horária mínima, não indica componentes curriculares específicos e não define o perfil profissional responsável pelo ensino. Em vez disso, atribui aos Estados, aos Municípios e ao Distrito Federal a definição dos parâmetros e das abordagens pedagógicas de implementação \cite{brasil2022resolucao1}.

Essa abertura permite adaptações às condições locais. Ao mesmo tempo, transfere para as redes e instituições o trabalho de transformar orientações gerais em decisões concretas, ou seja, definir como os conhecimentos serão distribuídos, quem ficará responsável por eles, que formação será oferecida e quais recursos serão necessários.

\subsection{A política nacional de educação digital}
\label{subsec:politica-nacional-educacao-digital}

O ambiente normativo continuou a se modificar depois da aprovação da BNCC Computação. A Lei nº 14.533/2023 instituiu a \sigla{PNED}{Política Nacional de Educação Digital}, cujo escopo inclui educação digital escolar, capacitação e especialização digital, pesquisa e desenvolvimento e inclusão digital \cite{brasil2023pned}.

No campo curricular, a PNED abrange uma agenda mais ampla do que a BNCC Computação. Além dos conhecimentos próprios da Computação, envolve letramento digital, educação midiática e uso pedagógico de dispositivos e recursos tecnológicos \cite{brasil2023pned}.

A Resolução CNE/CEB nº 2/2025 apresentou diretrizes para o uso de dispositivos digitais e para a integração curricular da educação digital e midiática. O documento retomou os eixos da BNCC Computação e explicitou a possibilidade de implementação disciplinar, transversal ou combinada, com implementação obrigatória a partir de 2026 \cite{brasil2025resolucao2}.

A BNCC Computação e a PNED podem ser compreendidas de forma articulada, mas não são equivalentes. A primeira define aprendizagens específicas de Computação. A segunda trata de uma política mais ampla de educação digital. Essa distinção revela que o ensino dos conhecimentos próprios da Computação está sendo acompanhado por uma agenda de uso de recursos tecnológicos.

\subsection{Formação docente: exigências e indefinições normativas}
\label{subsec:formacao-docente-normas}

A Resolução CNE/CEB nº 1/2022 determina que a formação inicial e continuada considere as aprendizagens da BNCC Computação \cite{brasil2022resolucao1}. Para a formação continuada, a Resolução CNE/CEB nº 2/2025 estabelece que cada rede deve elaborar um plano de formação coerente com a estratégia escolhida para implementar a educação digital \cite{brasil2025resolucao2}.

Quanto à formação inicial, entretanto, a orientação permanece genérica. A Resolução CNE/CEB nº 1/2022 inclui essa formação em seu campo de aplicação, mas não detalha como a exigência deve ser incorporada aos currículos das licenciaturas \cite{brasil2022resolucao1}. A Resolução CNE/CEB nº 2/2025 concentra-se na formação continuada dos profissionais que atuam na Educação Básica e não acrescenta orientações específicas para a organização curricular da formação inicial \cite{brasil2025resolucao2}. Assim, as normas não definem quais licenciaturas devem incorporar conhecimentos de Computação, quais conhecimentos devem integrar esses cursos, com que profundidade devem ser abordados, em quais núcleos formativos devem ser inseridos ou qual carga horária deve ser destinada a eles.

Desse modo, as normas apresentam maior detalhamento sobre as aprendizagens que devem ser garantidas aos estudantes da Educação Básica e sobre a formação continuada necessária à implementação do que sobre as mudanças curriculares requeridas na formação inicial de professores.

Na prática, essa ausência de detalhamento impõe que as instituições formadoras interpretem a exigência normativa e decidam como os currículos de seus cursos de licenciatura responderão a ela.

\section{Desafios de implementação do ensino de Computação}

\subsection{Desafios interdependentes}
\label{subsec:desafios-interdependentes}

A existência de um currículo prescrito não garante que ele seja realizado na prática da mesma forma como foi formulado. Entre a definição das orientações curriculares e o ensino em sala de aula, os objetivos e os conteúdos passam por diferentes níveis de interpretação e decisão. Nesse percurso, as orientações são transformadas em currículos locais, materiais, planejamentos, atividades e formas de avaliação \cite{remillard2014conceptualizing}. No ensino de Computação, esse processo também depende da maneira como os professores compreendem o currículo e incorporam suas orientações às escolhas de conteúdos, às estratégias de ensino e à avaliação da aprendizagem \cite{yiatas2026enacting}.

A literatura organiza os fatores envolvidos nesse processo de diferentes maneiras. \citeonline{remillard2014conceptualizing} analisam a passagem entre a formulação e a realização do currículo. O framework \sigla{EnCITE}{Entrypoints to Computing Integrated Teacher Education}, elaborado para discutir a integração da Computação e das tecnologias digitais à formação de professores, distingue problemas e possibilidades de natureza tecnológica, pedagógica, ideológica, política e de desenvolvimento profissional \cite{vogel2024encite}. No contexto da implementação escolar da Computação, o relatório \textit{After the reboot} examina conjuntamente questões relacionadas ao currículo, à formação e à disponibilidade de professores, aos recursos e às condições oferecidas pelas escolas \cite{royal2017afterreboot}.

As classificações dessas referências possuem finalidades e níveis de análise diferentes e não correspondem diretamente umas às outras. Para organizar a discussão desenvolvida neste capítulo, os fatores mais relacionados ao problema investigado foram reunidos em quatro dimensões: curricular, formativa, infraestrutural e institucional. Esse agrupamento constitui uma síntese analítica elaborada pelo autor, e não uma taxonomia consolidada na literatura. Ele também considera os resultados de um estudo preliminar sobre a implementação da BNCC Computação em redes públicas paulistas, coautorado pelo autor desta pesquisa \cite{santana2026desafios}.

A primeira dimensão é \textbf{curricular}. A publicação da BNCC Computação define aprendizagens que devem ser consideradas pelas redes, mas não determina integralmente como elas serão organizadas e ensinadas em cada contexto. Redes e escolas ainda precisam definir objetivos, conteúdos, progressão, carga horária, formas de oferta e critérios de avaliação. Essas decisões fazem parte da passagem entre o currículo formal e o currículo realizado \cite{remillard2014conceptualizing}. Quando as definições permanecem vagas, os professores podem ter dificuldade para compreender o que deve ser ensinado, em que sequência os conhecimentos devem aparecer e como a aprendizagem será avaliada. Essa dificuldade também foi identificada no estudo preliminar coautorado pelo autor \cite{santana2026desafios}.

A segunda dimensão é \textbf{formativa}. Para interpretar o currículo e transformá-lo em atividades de ensino, o professor precisa conhecer os conteúdos da área, compreender sua progressão ao longo da Educação Básica, selecionar exemplos adequados e reconhecer dificuldades dos estudantes. Professores que não tiveram contato com esses conhecimentos durante sua formação podem não dispor do repertório necessário para tomar essas decisões. Estudos sobre desenvolvimento profissional em Computação indicam que a formação pode ampliar o conhecimento e a confiança dos professores, mas seus efeitos variam de acordo com a duração das ações, sua relação com a prática e o apoio oferecido depois da formação \cite{ma2025effectiveness,liu2024bringing,mouza2022investigating}. A importância de articular conhecimentos, pedagogia e desenvolvimento profissional também aparece entre as dimensões consideradas pelo EnCITE \cite{vogel2024encite}.

A terceira dimensão é \textbf{infraestrutural}. O ensino de Computação não depende exclusivamente de equipamentos, pois parte das aprendizagens pode ser desenvolvida por meio de atividades desplugadas. Ainda assim, dispositivos, conectividade, manutenção, espaços adequados, programas e materiais didáticos influenciam o que pode ser realizado. Esses recursos tornam-se especialmente importantes quando os objetivos envolvem programação, produção de artefatos digitais, comunicação em rede ou observação prática do funcionamento de sistemas computacionais. Relatórios sobre a implementação da Computação nas escolas mostram que limitações de infraestrutura podem restringir as atividades disponíveis aos professores, mesmo quando existe um currículo definido e há interesse em ensiná-lo \cite{royal2017afterreboot}. Dificuldades relacionadas à disponibilidade e às condições de uso dos recursos também foram encontradas no estudo preliminar brasileiro \cite{santana2026desafios}.

A quarta dimensão é \textbf{institucional}. A implementação não depende apenas das decisões tomadas individualmente pelos professores. Redes e escolas precisam definir quem será responsável pelo ensino, como os diferentes componentes serão coordenados, que tempo será destinado ao planejamento e à formação e como as ações serão acompanhadas. Também precisam criar condições para que a mudança não dependa apenas da iniciativa de poucos profissionais. A literatura sobre mudança educacional distingue a adoção formal de uma proposta de sua institucionalização, que ocorre quando ela passa a integrar as responsabilidades, os procedimentos e as rotinas da organização \cite{fullan1991meaning}. No campo da formação de professores, \citeonline{kellerfranco2018formacao} alertam que uma proposta pode adotar uma linguagem de inovação e, ainda assim, conservar a estrutura curricular e as práticas anteriores. Aspectos políticos, organizacionais e de distribuição de responsabilidades também são considerados pelo EnCITE \cite{vogel2024encite}. No estudo preliminar coautorado pelo autor, a falta de definição das responsabilidades e de articulação entre os participantes apareceu como uma dificuldade para transformar a BNCC Computação em ações contínuas \cite{santana2026desafios}.

As quatro dimensões se relacionam e podem limitar umas às outras. Um currículo bem definido pode não ser realizado quando os professores não receberam preparação para ensiná-lo. Uma formação pode produzir novos conhecimentos, mas ter pouco efeito quando não há tempo para planejamento, recursos ou apoio da gestão. Equipamentos podem permanecer subutilizados quando não existe clareza sobre os objetivos curriculares. Da mesma forma, ações iniciadas por professores ou projetos específicos podem desaparecer quando não são incorporadas às responsabilidades e às rotinas da instituição.

Esses problemas não aparecem como obstáculos independentes. Dessa forma, a introdução de um currículo de Computação exige ações articuladas sobre o que será ensinado, quem realizará o ensino, quais recursos estarão disponíveis e como a implementação será sustentada institucionalmente.

\subsection{Desafios específicos do ensino de Computação}

Parte dos obstáculos descritos também ocorre em outras reformas educacionais. Entretanto, a implementação do ensino de Computação apresenta dificuldades próprias, relacionadas à novidade da área nos currículos, à indefinição sobre os profissionais responsáveis e à proximidade aparente entre ensinar Computação e utilizar tecnologias digitais.

Uma dificuldade decorre do estágio ainda recente de institucionalização da área nos currículos escolares. Diferentemente de componentes consolidados, a Computação ainda não possui, em muitas redes, uma tradição curricular, um corpo docente claramente responsável ou uma organização de conteúdos já incorporada às rotinas das escolas \cite{royal2017afterreboot,palop2025redefining}.

No Brasil, estudos realizados depois da aprovação da BNCC Computação mostram que essa incorporação ainda está em curso. Em uma pesquisa com 213 professores da Educação Infantil de uma rede pública paulista, o conhecimento sobre a norma era baixo e suas habilidades não eram consideradas suficientemente claras pelos participantes \cite{guarda2024bnccinfantil}. Em escolas estaduais de dois municípios pernambucanos, 63\% das unidades investigadas ainda não ensinavam Computação, e os participantes apontaram dificuldades relacionadas à infraestrutura, à conectividade e à falta de professores qualificados \cite{oliveira2024ensinocomputacao}. Em outra investigação, realizada com dirigentes municipais de educação do Sertão do Submédio São Francisco, parte das redes ainda não conhecia suficientemente as normas, e a formação dos professores e as condições de infraestrutura apareciam como obstáculos à implementação \cite{souza2024implantacao}.

Esses resultados se referem a redes e localidades específicas e não permitem estimar a situação de todo o país. Ainda assim, mostram casos recentes em que a Computação não estava plenamente incorporada aos currículos, às responsabilidades docentes e às rotinas escolares.

Países com currículos de Computação mais antigos também enfrentam dificuldades relacionadas à disponibilidade de professores. Nos Estados Unidos, a escassez de docentes preparados para ensinar Ciência da Computação é apontada como uma das principais barreiras à ampliação do acesso à área \cite{hardin2026spark}.

Outra dificuldade é a proximidade aparente entre ensino de computação e uso de tecnologia. Como escolas e professores já utilizam plataformas, aplicativos e dispositivos, a implementação da BNCC Computação pode ser interpretada como ampliação dessas práticas. Essa interpretação encobre a diferença estabelecida na Seção~\ref{sec:computacao-conhecimento-escolar}: uma tecnologia pode mediar uma aula sem se tornar objeto de aprendizagem \cite{royal2017afterreboot,sadik2024understanding}.

A flexibilidade das formas de oferta também dificulta a implementação. A computação pode ser organizada como disciplina, distribuída transversalmente entre componentes ou implementada por uma combinação dessas estratégias \cite{brasil2025resolucao2}. Essa flexibilidade favorece a adaptação, mas amplia o número de decisões necessárias. Em uma oferta disciplinar, é preciso garantir espaço no currículo e profissionais responsáveis. Em uma oferta transversal, é preciso coordenar diferentes componentes para evitar fragmentação, lacunas e repetições superficiais.

Soma-se a isso a disponibilidade limitada de professores com formação específica em Computação. Em 2022, 257.581 estudantes concluíram cursos de licenciatura no Brasil, considerando as modalidades presencial e a distância \cite{inep2024resumotecnico2022}. No mesmo ano, o relatório da Sociedade Brasileira de Computação registrou 418 concluintes na categoria ``Licenciatura em Computação'' \cite{sbc2022estatisticas}\footnote{Os dois valores se referem a 2022, mas não permitem calcular a participação exata da Licenciatura em Computação no total das licenciaturas. O dado do \sigla{Inep}{Instituto Nacional de Estudos e Pesquisas Educacionais Anísio Teixeira} reúne os concluintes de cursos presenciais e a distância. O relatório da SBC não desagrega os 418 concluintes por modalidade de ensino e informa que cursos de bacharelado ou licenciatura da área de Computação com denominações diferentes das previstas na Resolução CNE/CES nº 5/2016 foram reunidos na categoria ``Outros Cursos''. Por essas razões, os valores são utilizados apenas para comparar a escala dos fluxos anuais, e não para produzir uma proporção exata.} O número mostra que o fluxo anual de concluintes registrado nessa categoria é reduzido diante do conjunto das licenciaturas.

Tal escassez não é exclusiva do sistema educacional brasileiro. No Reino Unido, dados oficiais do governo referentes a 2016 mostram que apenas 30\% dos professores de Tecnologia da Informação e Comunicação possuíam um diploma relevante na área, proporção inferior à observada entre professores de física (51\%) e de matemática (46\%) \cite{royal2017afterreboot}. Uma revisão internacional sobre a formação de professores de Ciência da Computação em diferentes países mostra que essa dependência de docentes de outras áreas também se repete fora do Reino Unido \cite{yadav2022review}. No Texas, por exemplo, mais de 80\% dos professores certificados para ensinar Computação obtiveram essa certificação depois de já estarem certificados em outra disciplina \cite{yadav2022review}.

\subsection{Evidências do contexto brasileiro}
\label{subsec:evidencias-contexto-brasileiro}

No Brasil, a implementação do ensino de Computação ocorre em um contexto marcado por desigualdades nas condições materiais e formativas das escolas. A pesquisa TIC Educação, realizada pelo \sigla{Cetic.br}{Centro Regional de Estudos para o Desenvolvimento da Sociedade da Informação} desde 2010, mostra que a existência de conexão à Internet não significa, necessariamente, que computadores e outros dispositivos estejam disponíveis aos estudantes ou possam ser utilizados regularmente nas atividades escolares \cite{cetic2024tic2023}. Portanto, a infraestrutura não pode ser analisada apenas pela presença de conectividade. Também devem ser consideradas a disponibilidade efetiva de equipamentos, sua manutenção, as condições de acesso e as possibilidades de uso pedagógico.

A extensão da lacuna formativa é mais difícil de determinar. O Indicador de Adequação da Formação Docente, produzido pelo Inep, classifica cada docência de acordo com a relação entre a formação inicial do professor e a disciplina que ele leciona. Para realizar essa classificação, a metodologia do indicador estabelece correspondências entre os componentes curriculares da Educação Básica e as áreas de formação superior consideradas adequadas para seu ensino \cite{inep2014adequacaoformacaodocente,inep2021atualizacaoadequacao}. Na tabela utilizada para estabelecer essas correspondências, a Computação não aparece como um componente curricular analisado separadamente. Desse modo, os resultados do indicador não permitem identificar quantos professores possuem formação considerada adequada especificamente para ensinar os conhecimentos previstos na BNCC Computação. Essa ausência limita tanto o diagnóstico da situação atual quanto o planejamento e o acompanhamento das políticas de formação. A própria Sociedade Brasileira de Computação recomenda a criação de indicadores nacionais para monitorar os efeitos da inserção da Computação nas escolas \cite{sbc2025grandesdesafios}.

Em seu documento sobre os grandes desafios da Educação em Computação para o período de 2025 a 2035, a Sociedade Brasileira de Computação caracteriza a formação docente e a produção de recursos pedagógicos como um dos principais desafios nacionais da área. O documento estima que apenas entre 10 e 15 mil, dentre mais de 2 milhões de professores brasileiros, possuam formação em Computação, o que expressa a magnitude do desequilíbrio identificado pela comunidade brasileira de Educação em Computação entre a quantidade de professores da Educação Básica e a disponibilidade de profissionais com formação específica \cite{sbc2025grandesdesafios}.

O documento também identifica fatores que dificultam a superação desse desequilíbrio. Há poucos cursos de Licenciatura em Computação: o Cadastro e-MEC do Ministério da Educação registrava, em julho de 2026, apenas 94 cursos ativos com essa denominação em todo o país, dos quais 66 já tinham funcionamento iniciado\footnote{Levantamento no Cadastro e-MEC (filtros: Graduação, curso Computação, grau Licenciatura, situação Em Atividade), com 97 registros brutos; excluídos 3 cursos sem relação direta com a formação de professores de Computação, restando 94. O e-MEC não gera link estável para consultas filtradas, por isso o relatório foi preservado pelo autor.} \cite{brasil2026emec}. Esse número é pequeno diante da rede que precisará desses professores: o Censo Escolar de 2025 registrou 178,8 mil escolas de educação básica e 2,4 milhões de professores em atividade no país \cite{inep2025apresentacaocenso2024}. Além disso, profissionais formados na área podem optar por outras oportunidades de trabalho, enquanto professores que já atuam nas escolas enfrentam sobrecarga, condições de trabalho desfavoráveis e dificuldades para participar de processos contínuos de formação \cite{sbc2025grandesdesafios}. Esses fatores indicam que a expansão do ensino de Computação dificilmente poderá depender apenas da disponibilidade de licenciados na área.

Estudos com egressos de cursos de Licenciatura em Computação mostram que parte deles não segue a docência da área na Educação Básica. Um levantamento nacional com 60 egressos das cinco regiões brasileiras mostra que 51,67\% declararam atuar no ensino de Computação, enquanto 48,33\% informaram não atuar nessa atividade \cite{teixeira2019perfil}. O mesmo levantamento mostra que 88,3\% desses egressos avaliaram haver dificuldade ou muita dificuldade para a inserção profissional \cite{teixeira2019perfil}.

Estudos realizados em cursos específicos mostram proporções de inserção na docência ainda menores. Uma pesquisa com 22 egressos de uma universidade pública de Mato Grosso identificou apenas 2 participantes atuando na docência da Educação Básica \cite{oliveira2018insercao}. Um estudo com egressos da Universidade de Brasília mostra que 12 entre 18 participantes de grupos focais trabalhavam em atividades técnicas de Tecnologia da Informação \cite{pinheiro2017egressos}. Entre esses 12, apenas 3 exerciam atividades relacionadas diretamente à Computação e à Educação \cite{pinheiro2017egressos}. Um levantamento com egressos do Instituto Federal do Tocantins mostra que, entre 59 respondentes, 15 atuavam na docência \cite{rodrigues2020perfil}. Desses 15, 30\% ministravam disciplinas relacionadas à informática e 70\% ensinavam outras disciplinas \cite{rodrigues2020perfil}.

Essas evidências são predominantemente amostrais e locais. Os estudos também definem de formas distintas o que conta como atuação no ensino de Computação. Essa diferença de critérios impede a agregação direta dos resultados em uma estimativa nacional. A hipótese de que os egressos migram para atividades técnicas por causa da remuneração não é testada em nenhum desses estudos.

A insuficiência formativa também possui uma dimensão qualitativa. Segundo o documento da Sociedade Brasileira de Computação (2025), parte dos gestores e professores ainda confunde o ensino de Computação com o uso de ferramentas digitais. Essa compreensão restringe a implementação a atividades com plataformas, aplicativos ou dispositivos, sem assegurar a aprendizagem de conceitos relacionados ao pensamento computacional, ao mundo digital e à cultura digital \cite{sbc2025grandesdesafios}. O problema, portanto, não consiste somente na ausência de profissionais formalmente habilitados, mas também na falta de conhecimentos que permitam reconhecer a Computação como objeto de ensino, interpretar as habilidades curriculares e transformá-las em objetivos, atividades e formas de avaliação.

A Sociedade Brasileira de Computação também relaciona a insuficiência da formação aos riscos de um ensino descontextualizado, sem objetivos pedagógicos claros, e à possível intensificação da sobrecarga física e emocional dos professores. Ao mesmo tempo, destaca a carência de materiais didáticos acessíveis e adequados à diversidade cultural e social brasileira \cite{sbc2025grandesdesafios}. Desse modo, a preparação docente não pode ser separada das condições de trabalho, da disponibilidade de recursos, da definição das responsabilidades e do apoio oferecido pelas redes e instituições.

Entre as ações recomendadas, o documento propõe ampliar os cursos de Licenciatura em Computação, inserir conhecimentos da área em cursos de Pedagogia e em licenciaturas de outros campos, estruturar currículos de referência, fortalecer a formação continuada e estabelecer parcerias entre universidades, escolas e governos \cite{sbc2025grandesdesafios}. A inclusão das demais licenciaturas nessa agenda é particularmente relevante porque reconhece que professores com diferentes formações poderão participar da implementação da Educação em Computação. Entretanto, a recomendação permanece em nível geral: não determina quais licenciaturas devem ser modificadas, quais conhecimentos devem ser incorporados, com que profundidade, em quais espaços curriculares ou de que maneira essas mudanças podem ser compatibilizadas com as condições de cada curso.

Estudos brasileiros sobre a implementação da BNCC Computação apresentam dificuldades convergentes, relacionadas à formação docente, à infraestrutura, à interpretação das habilidades, à disponibilidade de materiais e à articulação institucional. O estudo realizado pelo grupo de pesquisa ao qual este trabalho está vinculado também identificou que a principal dificuldade não está apenas em conhecer a norma, mas em traduzi-la para a organização curricular e para a rotina escolar \cite{santana2026desafios}.

Os resultados desse estudo mostram que parte das ações interpretadas como implementação da BNCC Computação permanece concentrada no uso de plataformas e ferramentas digitais, sem garantir o ensino dos conhecimentos previstos nos três eixos. Também indicam que formação, infraestrutura, currículo e governança precisam ser tratados de maneira articulada. Nesta seção, esses resultados são utilizados como evidência do problema de implementação. O delineamento, os participantes, os resultados completos e as limitações do estudo são apresentados na Seção~\ref{sec:estudo-implementacao-bncc}.

Em conjunto, essas evidências mostram que a formação docente ocupa uma posição central, mas não isolada, na implementação do ensino de Computação. É por meio dela que os professores podem compreender a especificidade da área, desenvolver conhecimentos sobre seus conteúdos e construir formas de ensiná-los. Entretanto, seus efeitos dependem de condições materiais, recursos pedagógicos, tempo, acompanhamento e apoio institucional.

\section{A formação docente e a importância da formação inicial}
\label{sec:formacao-inicial-estrutural}

Ensinar Computação exige mais do que saber utilizar uma ferramenta digital. O professor precisa compreender os conhecimentos previstos no currículo, relacioná-los à etapa de escolarização, selecionar formas adequadas de apresentá-los e acompanhar a aprendizagem dos estudantes \cite{shulman1986understanding,shulman1987knowledge}. Também precisa reconhecer quando uma atividade tem como objetivo desenvolver conhecimentos próprios da Computação e quando apenas utiliza um recurso digital como meio para ensinar outro conteúdo, distinção discutida na Seção~\ref{sec:computacao-conhecimento-escolar} \cite{royal2012shut}.

A formação docente é, portanto, uma condição para transformar as aprendizagens prescritas em experiências de ensino. Pesquisas sobre desenvolvimento profissional mostram efeitos positivos sobre o conhecimento e a prática de professores de Computação \cite{ma2025effectiveness}. Estudos também registram ganhos de conhecimento conceitual e de autoeficácia depois de formações mais longas e relacionadas à prática \cite{nugent2022developing,liu2025building}.

Entretanto, os resultados das formações não dependem apenas de seu conteúdo. Professores que desenvolvem novos conhecimentos podem encontrar dificuldades para utilizá-los quando não há tempo, apoio da gestão, recursos ou articulação com o currículo escolar. O acompanhamento realizado por \citeonline{mouza2022investigating} mostrou que professores submetidos a uma mesma formação modificaram suas práticas de maneiras diferentes, em razão das condições encontradas nas escolas em que atuavam. Esse resultado indica que a preparação individual do professor, embora necessária, não é suficiente para produzir mudanças na prática.

Por essa razão, a formação deve ser compreendida como um dos componentes de um sistema mais amplo de implementação. Ela pode preparar o professor para compreender os conteúdos e desenvolver formas de ensiná-los, mas seus efeitos dependem da articulação com decisões curriculares, infraestrutura, apoio institucional e definição de responsabilidades.

\subsection{Limites da formação continuada isolada}
\label{subsec:limites-formacao-continuada}

Para os professores já em exercício, a formação continuada é a resposta mais imediata às novas demandas. Ela pode apresentar conteúdos, atualizar conhecimentos, apoiar o planejamento e acompanhar mudanças curriculares. Sua importância aumenta em uma área que se modifica rapidamente e cuja presença na formação inicial ainda é limitada.

A literatura indica que muitas iniciativas de formação continuada em Pensamento Computacional ou programação são breves, o que limita as oportunidades de aprofundamento e de acompanhamento de seus efeitos na prática docente. Uma revisão internacional identificou o predomínio de programas realizados em poucos dias, enquanto iniciativas desenvolvidas ao longo de um ano letivo eram minoritárias \cite{liu2024bringing}. Nesse conjunto de estudos, os resultados observados concentravam-se principalmente na familiarização com conceitos, no aumento da confiança dos professores e na elaboração inicial de planos de aula. Em contrapartida, eram menos frequentes as evidências de mudanças sustentadas nas práticas de ensino e de efeitos sobre a aprendizagem dos estudantes.

No Brasil, um mapeamento de unidades instrucionais voltadas ao ensino de algoritmos e programação também encontrou predomínio de oficinas curtas e concentração em conteúdos técnicos \cite{kretzer2020formacao}. Outra revisão identificou que menos da metade dos estudos brasileiros analisados avaliava efetivamente as habilidades desenvolvidas pelos professores \cite{sokolonski2021revisao}.

Formações curtas podem cumprir objetivos introdutórios e atender a necessidades imediatas. O problema aparece quando precisam compensar, sozinhas, a ausência de bases conceituais e pedagógicas que poderiam ter sido construídas ao longo da formação profissional.

A continuidade das iniciativas também pode depender de financiamento temporário, equipes externas e adesão voluntária de professores já interessados no tema \cite{basu2021building}. Quando essas condições desaparecem, a formação pode não alcançar novos profissionais nem produzir mudanças duradouras nas instituições.

A formação continuada permanece necessária para professores em exercício e para a atualização de quem já recebeu uma preparação inicial. Entretanto, seus limites indicam que a resposta à BNCC Computação não pode depender exclusivamente de cursos pontuais oferecidos depois da entrada na profissão.

\subsection{Possibilidades estruturais da formação inicial}
\label{subsec:possibilidades-formacao-inicial}

A formação inicial oferece condições diferentes das encontradas em cursos e oficinas de formação continuada. Uma licenciatura se desenvolve durante vários semestres e reúne componentes dedicados aos conhecimentos específicos da área, à formação pedagógica, às práticas de ensino e ao estágio supervisionado \cite{brasil2024resolucao4}. Essa organização permite que um novo conhecimento seja trabalhado em mais de um momento do curso e relacionado a diferentes partes da preparação profissional, em vez de ser apresentado apenas em uma atividade isolada.

Essa possibilidade é importante porque aprender a ensinar um conteúdo exige tempo e diferentes tipos de experiência. O licenciando precisa, primeiro, compreender os conceitos da área. Depois, precisa aprender a selecionar o que é adequado para cada etapa escolar, pensar em exemplos e atividades, antecipar dificuldades dos estudantes e avaliar o que foi aprendido. A literatura sobre formação de professores mostra que essa preparação se torna mais consistente quando os componentes do curso se relacionam entre si e quando os conhecimentos estudados são retomados em situações de planejamento, prática e reflexão \cite{DarlingHammond2005Design,hammerness2006coherence}.

No caso da Computação, isso significa que o primeiro contato com seus conhecimentos pode ocorrer em um componente introdutório e ser retomado posteriormente em disciplinas de didática, práticas de ensino, projetos ou estágios. Nessas retomadas, o foco pode mudar. Em um primeiro momento, o licenciando pode estudar conceitos como algoritmos, representação de dados ou funcionamento de sistemas digitais. Em outro, pode analisar como esses conhecimentos aparecem na BNCC Computação, como podem ser ensinados a diferentes públicos e quais dificuldades podem surgir durante a aprendizagem. Estudos sobre a integração da Computação em programas de formação inicial mostram que atividades da área podem ser inseridas em diferentes componentes do curso, desde que essa distribuição seja planejada de acordo com os objetivos e com a organização de cada programa formativo \cite{margulieux2022integrating}.

Essa vantagem estrutural, entretanto, não garante que a formação seja adequada. Um curso pode incluir conteúdos de Computação de maneira superficial, concentrá-los em uma única disciplina sem relação com o restante da licenciatura ou separar o estudo dos conceitos das experiências de ensino. A literatura sobre formação docente também mostra que cursos longos podem permanecer fragmentados quando seus componentes não compartilham objetivos e não estabelecem relações claras entre teoria e prática \cite{hammerness2006coherence}.

Com base nessa literatura, este trabalho pressupõe que a formação inicial, por sua duração e por reunir diferentes espaços de formação, oferece condições para que os conhecimentos de Computação sejam introduzidos, retomados, aprofundados e relacionados à preparação pedagógica. Para que essa possibilidade se concretize, é necessário decidir quais conhecimentos serão incluídos, em que momentos do curso aparecerão e como serão articulados às finalidades e às condições de cada licenciatura.

\subsection{Definição de integração curricular adotada}
\label{subsec:definicao-integracao-curricular}

\citeonline{burke2023conceptualising} mostram que a integração curricular pode assumir diferentes formas. Para compreender o que é uma proposta de integração e como ela é organizada, os autores sugerem analisar três aspectos. O primeiro é a finalidade: por que determinados conhecimentos devem ser integrados e que aprendizagem se pretende promover. O segundo diz respeito aos conhecimentos envolvidos: quais conteúdos, áreas ou experiências serão relacionados e como essa relação será construída. O terceiro considera o modo como a integração responde às pessoas e ao contexto em que ocorre, incluindo os estudantes, os professores, as orientações curriculares e as condições de realização.

Esses aspectos mostram que colocar conteúdos diferentes em uma mesma disciplina, atividade ou tema não garante, por si só, que exista integração curricular. É necessário explicar por que esses conhecimentos estão sendo aproximados, qual papel cada um desempenha na formação e de que maneira serão relacionados durante o curso. Sem essas definições, a aproximação pode permanecer apenas formal, sem produzir uma formação articulada.

Nesta pesquisa, a expressão \textit{integração curricular da Computação} designa o processo de incorporar conhecimentos e práticas da Computação ao currículo de uma licenciatura de forma planejada e relacionada às finalidades do curso. Isso exige considerar o perfil profissional que a licenciatura pretende formar, as responsabilidades que seus egressos poderão assumir na Educação Básica e os conhecimentos de Computação necessários para essa atuação.

Com base nessa definição, a integração curricular é compreendida como um conjunto de decisões interdependentes, e não apenas como a escolha de um componente no qual um novo conteúdo será inserido. \citeonline{tyler1949principles} organizou o planejamento curricular em torno da definição dos objetivos educacionais, da seleção das experiências de aprendizagem, da organização dessas experiências e da avaliação dos resultados alcançados. Desenvolvimentos posteriores da teoria curricular também destacam que as decisões sobre finalidades, conteúdos, organização, experiências formativas e avaliação precisam ser consideradas de forma articulada e relacionadas ao contexto em que o currículo será realizado \cite{ornsteinhunkins2009curriculum,pacheco2001curriculo,gasparroldao2007elementos}.

Nesta pesquisa, essas contribuições são utilizadas para explicitar seis dimensões da integração curricular da Computação: a finalidade da formação; os conhecimentos a serem contemplados; a profundidade com que serão desenvolvidos; sua progressão ao longo do curso; os espaços curriculares utilizados; e as práticas formativas e formas de avaliação. O eixo relativo às pessoas e ao contexto, presente na formulação de \citeonline{burke2023conceptualising}, é considerado por meio da análise dos atores envolvidos, das condições do curso, das restrições normativas e institucionais e da possibilidade de executar e sustentar a mudança.

A integração curricular pode ampliar o contato dos licenciandos com a Computação, mas não assegura que os conhecimentos estejam relacionados ao perfil do egresso, distribuídos de maneira coerente ao longo da formação ou articulados à preparação pedagógica. Estudos sobre a inserção da Computação na formação inicial também mostram que as mudanças precisam ser definidas a partir das necessidades de cada programa e relacionadas aos componentes nos quais serão desenvolvidas \cite{margulieux2022integrating}.

\subsection{Conhecimentos necessários à preparação docente}
\label{subsec:conhecimentos-preparacao-docente}

A preparação docente não pode ser reduzida à aquisição de domínio técnico. Em seus trabalhos sobre o conhecimento profissional docente, \citeauthoronline{shulman1986understanding} (\citeyear{shulman1986understanding}, \citeyear{shulman1987knowledge}) discute a separação rígida entre domínio do conteúdo e conhecimento pedagógico. Para ensinar, o professor precisa compreender o conteúdo e saber transformá-lo em formas acessíveis aos estudantes.

Aplicada à preparação de professores para o ensino de Computação, essa distinção permite identificar pelo menos três planos relacionados. O primeiro é o domínio técnico, que envolve operar ferramentas, linguagens, dispositivos ou ambientes. O segundo é o conhecimento do conteúdo, relacionado à compreensão dos conceitos, das práticas e dos problemas da Computação que integram a formação prevista para a Educação Básica. O terceiro é o conhecimento pedagógico do conteúdo, que envolve selecionar representações, antecipar dificuldades, organizar progressões, propor atividades e avaliar a aprendizagem \cite{saeli2011teaching,yadav2019computer}.

Esses planos se relacionam, mas não são equivalentes. Uma formação centrada no uso de uma ferramenta pode desenvolver habilidades operacionais sem produzir compreensão conceitual. De modo semelhante, conhecer conceitos de Computação não garante que o futuro professor saiba ensiná-los a crianças ou adolescentes.

O conhecimento pedagógico do conteúdo também depende do público e da etapa escolar. Um exemplo, uma representação ou uma dificuldade comum nos Anos Iniciais pode não ser adequado ao Ensino Médio. A preparação precisa, portanto, relacionar os conhecimentos da Computação aos estudantes, aos objetivos e aos contextos de atuação de cada licenciatura.

\subsection{Diretrizes brasileiras e complementaridade entre formações}
\label{subsec:diretrizes-complementaridade-formacao}

A Resolução CNE/CP nº 4/2024 redefiniu as Diretrizes Curriculares Nacionais para a formação inicial de profissionais do magistério \cite{brasil2024resolucao4}. O documento exige que as instituições formadoras promovam competências relacionadas às Tecnologias Digitais de Informação e Comunicação e inclui o domínio dessas tecnologias entre as capacidades esperadas dos egressos.

A resolução também determina que os núcleos de aprofundamento dos conteúdos específicos garantam a compreensão do conhecimento pedagógico do conteúdo necessário ao planejamento e à realização de situações de ensino e aprendizagem \cite{brasil2024resolucao4}. Essa organização oferece uma base para articular conteúdo e pedagogia, mas a norma não menciona diretamente a Computação ou o Pensamento Computacional.

A exigência permanece no nível mais amplo das competências digitais. Não se define quais licenciaturas devem tratar dos conhecimentos próprios da Computação, quais objetivos devem ser alcançados ou em quais núcleos e componentes esses conhecimentos devem ser inseridos.

Reconhecer a importância da formação inicial não implica atribuir a ela toda a solução. Mesmo a Licenciatura em Computação apresenta tensões entre o perfil de formação, a atuação profissional e o espaço destinado ao ensino de Computação na Educação Básica \cite{gentil2013licenciatura}. Além disso, currículos possuem limites de carga horária, o número de formadores preparados é reduzido e os conhecimentos e tecnologias da área continuam se modificando.

A formação inicial e a continuada devem ser compreendidas como complementares. A formação inicial pode construir bases conceituais e pedagógicas de forma progressiva. A formação continuada pode atualizar, aprofundar e apoiar a implementação em contextos específicos. A questão que permanece é como essa preparação inicial deve variar entre cursos que formam profissionais com responsabilidades distintas na Educação Básica.

\section{Considerações finais}
\label{sec:consideracoes-finais-referencial}

Este capítulo mostrou que a implementação da BNCC Computação depende da articulação entre decisões curriculares, formação docente, infraestrutura e apoio institucional. Nesse contexto, a formação continuada é necessária, mas não deve constituir a única resposta, sobretudo quando precisa compensar, por meio de ações pontuais, conhecimentos conceituais e pedagógicos que não foram desenvolvidos anteriormente.

A formação inicial oferece condições para que esses conhecimentos sejam trabalhados de forma progressiva e relacionados à preparação pedagógica dos futuros professores. Essa formação, entretanto, não pode ser igual em todas as licenciaturas, pois os cursos possuem objetivos, currículos e perfis profissionais distintos. Torna-se necessário, portanto, examinar como a Computação pode ser integrada a esses cursos de acordo com suas finalidades formativas e condições institucionais. Essa questão orienta a discussão apresentada no capítulo seguinte.